% TOP_PROBLEM: {"id": 2800704, "title": "10 Lectures and 42 Open Problems — The Deletion Channel", "category_id": 15, "category": {"id": 15, "name": "computer_science", "display_name": "Computer Science", "description": "Computational complexity, algorithms, and theoretical CS.", "slug": "computer-science", "order_index": 15, "created_at": "2026-07-31T15:26:25.671Z"}, "status": "partially_solved", "problem_number": "AMR-027-0704", "set_id": 13, "statusQualification": "The source says to draw a random binary input. Average-case and worst-case trace complexity, and the chosen error probabilities, are explicitly distinguished."}
% FIRST_POSTED: 2026-10-02
% ENTRY_KIND: statement_only
% UnsolvedMath ID 2800704; source catalogue code AMR-027-0704
% !TeX program = lualatex
% Definition and sources only; this file is not an LLM solution attempt.
\documentclass[11pt,a4paper]{article}
\usepackage[margin=25mm]{geometry}
\usepackage{fontspec}
\setmainfont{lmroman10-regular.otf}[BoldFont=lmroman10-bold.otf,
 ItalicFont=lmroman10-italic.otf,BoldItalicFont=lmroman10-bolditalic.otf]
\usepackage{amsmath,amssymb,mathtools}
\usepackage{unicode-math}
\setmathfont{latinmodern-math.otf}
\AtBeginDocument{\let\setminus\smallsetminus}
\usepackage{xurl,hyperref}
\hypersetup{colorlinks=true,urlcolor=blue}
\setcounter{secnumdepth}{0}
\setlength{\parindent}{0pt}
\setlength{\parskip}{5pt}
\setlength{\emergencystretch}{3em}
\begin{document}
\section*{10 Lectures and 42 Open Problems — The Deletion Channel}\label{2800704}
{\small UnsolvedMath ID 2800704 · AMR-027-0704 · Computer Science · AMR Open Problem Lists}
\subsection{Definitions and mathematical statement}
A trace of $x\in\{0,1\}^n$ is obtained by independently deleting each
bit with probability $1/2$ and retaining the order of the remaining
bits. Deleted positions are not reported. Repeated traces of the
same input use independent deletions.

For $0<\varepsilon<1/2$, let $D_{\rm av}(n,\varepsilon)$ be the
least $m$ for which a decoder, knowing $n$, reconstructs a uniformly
random $X\in\{0,1\}^n$ from $m$ traces with error probability at most
$\varepsilon$. Probability averages over $X$, deletions, and any
decoder randomness. This makes explicit the source's random-input
wording. Define $D_{\rm worst}$ separately by requiring the error bound
for every fixed input string.

Determine the asymptotic trace complexity, specifying the convention
used; the high-probability average-case target can be fixed as
$D_{\rm av}(n,n^{-1})$ for $n\ge3$. Bounds proved only for
$D_{\rm worst}$ must be identified as such rather than silently
equating the two quantities.

The second question concerns two known distinct strings $x,y$.
Let $d(x,y)$ be the least number of traces permitting a test whose
error is at most $1/3$ under each of the two hypotheses. Determine
\[
                  \max_{x\ne y\in\{0,1\}^n}d(x,y)
\]
and identify pairs realizing the hardest order of growth. These
questions concern information from repeated noisy observations, not
the capacity of a code used for one transmission. Computation time is
unrestricted in these sample-complexity definitions.
\subsection{Short English statement}
Determine trace-reconstruction sample complexity and the hardest pairs of strings to distinguish after independent deletions.
\subsection{Sources}
\begin{itemize}
\item[S1] ULAM.AI and the UnsolvedMath Contributors, supplied problems.json, record 2800704 (AMR-027-0704), AMR Open Problem Lists. Catalogue identity and source transcription; CC BY 4.0.
\url{https://huggingface.co/datasets/ulamai/UnsolvedMath}.
\item[S2] Bandeira - 42 Open Problems in Data Science (2016), Open Problem 7.4. Original source indexed by AMR-027-0704.
\url{https://afonsobandeira.wordpress.com/2015/11/27/10l42grouptesting/}.
\item[S3] A. S. Bandeira, Ten Lectures and Forty-Two Open Problems in the Mathematics of Data Science, corresponding numbered problem and surrounding definitions.
\url{https://people.math.ethz.ch/~abandeira/TenLecturesFortyTwoProblems.pdf}.
\end{itemize}
\end{document}
